\documentclass[table]{beamer}
%[]中可以使用draft、handout、screen、transparency、trancompress、compress等参数

%指定beamer的模式与主题
\mode<presentation>
{
  \usetheme{Madrid}
%\usetheme{Boadilla}
%\usecolortheme{default}
%\usecolortheme{orchid}
%\usecolortheme{whale}
%\usefonttheme{professionalfonts}
}

%\usetheme{Madrid}
%这里还可以选择别的主题：Bergen, Boadilla, Madrid, AnnArbor, CambridgeUS, Pittsburgh, Rochester, Warsaw, ...
%有导航栏的Antibes, JuanLesPins, Montpellier, ...
%有内容的Berkeley, PaloAlto, Goettingen, Marburg, Hannover, ...
%有最小导航栏的Berlin, Ilmenau, Dresden, Darmstadt, Frankfurt, Singapore, Szeged, ...
%有章和节表单的Copenhagen, Luebeck, Malmoe, Warsaw, ...

%\usecolortheme{default}
%设置内部颜色主题(这些主题一般改变block里的颜色);这个主题一般选择动物来命名
%这里还可以选择别的颜色主题,如默认的和有特别目的的颜色主题default,structure,sidebartab,全颜色主题albatross,beetle,crane,dove,fly,seagull,wolverine,beaver

%\usecolortheme{orchid}
%设置外部颜色主题(这些主题一般改变title里的颜色);这个主题一般选择植物来命名
%这里还可以选择别的颜色主题,如默认的和有特别目的的颜色主题lily,orchid,rose

%\usecolortheme{whale}
%设置字体主题;这个主题一般选择海洋动物来命名
%这里还可以选择别的颜色主题,如默认的和有特别目的的颜色主题whale,seahorse,dolphin

%\usefonttheme{professionalfonts}
%类似的还可以定义structurebold,structuresmallcapsserif,professionalfonts

% 控制 beamer 的风格,可以根据自己的爱好修改
%\usepackage{beamerthemesplit} %使用 split 风格
%\usepackage{beamerthemeshadow} %使用 shadow 风格
%\usepackage[width=2cm,dark,tab]{beamerthemesidebar}

%插入音标
%\usepackage{tipa}
%\AtBeginDocument{
  %\renewcommand\textipa{\fontencoding{T3}\selectfont}
%}
%\AtBeginDocument{
  %\renewcommand\textipa[2][r]{{\fontfamily{cm#1}\tipaencoding #2}}
%}
%\renewenvironment{IPA}[1][r]
 %{\fontfamily{cm#1}\tipaencoding}
 %{}

% 设定英文字体
%\usepackage{fontspec}
% Fix bugs for fontspec in TeXLive2015
\ifdefined\suppressfontnotfounderror
  \expandafter\let\csname xetex_suppressfontnotfounderror:D\endcsname
    \suppressfontnotfounderror
\else
  \expandafter\let\csname xetex_suppressfontnotfounderror:D\endcsname
    \luatexsuppressfontnotfounderror
\fi
\usepackage[no-math]{fontspec}
\setmainfont{Times New Roman}
\setsansfont{Arial}
\setmonofont{Courier New}

% 设定中文字体
\usepackage[BoldFont,SlantFont,CJKchecksingle,CJKnumber]{xeCJK}
%\setCJKmainfont[BoldFont={Adobe Heiti Std},ItalicFont={Adobe Kaiti Std}]{Adobe Song Std}
\setCJKmainfont[BoldFont={Adobe Heiti Std},ItalicFont={Adobe Kaiti Std}]{WenQuanYi Micro Hei}
\setCJKsansfont{Adobe Heiti Std}
\setCJKmonofont{Adobe Fangsong Std}
\punctstyle{hangmobanjiao}

\defaultfontfeatures{Mapping=tex-text}
\usepackage{xunicode}
\usepackage{xltxtra}

\XeTeXlinebreaklocale "zh"
\XeTeXlinebreakskip = 0pt plus 1pt minus 0.1pt

\usepackage{setspace}
\usepackage{colortbl,xcolor}
\usepackage{hyperref}
%\hypersetup{xetex,bookmarksnumbered=true,bookmarksopen=true,pdfborder=1,breaklinks,colorlinks,linkcolor=blue,filecolor=black,urlcolor=cyan,citecolor=green}
\hypersetup{xetex,bookmarksnumbered=true,bookmarksopen=true,pdfborder=1,breaklinks,colorlinks,linkcolor=cyan,filecolor=black,urlcolor=blue,citecolor=green}

% 插入图片
\usepackage{graphicx}
\graphicspath{{figures/}}
% 图文混排
%\usepackage{picins}
\usepackage{floatflt}

% 可能用到的包
\usepackage{amsmath,amssymb}
%插入多媒体
%\usepackage{media9}
%\usepackage{movie15}
\usepackage{multimedia}
\usepackage{multicol}
\usepackage{multirow}

% 定义一些自选的模板，包括背景、图标、导航条和页脚等，修改要慎重
% 设置背景渐变由10%的红变成10%的结构颜色
%\beamertemplateshadingbackground{red!10}{structure!10}
%\beamertemplatesolidbackgroundcolor{white!90!blue}
% 使所有隐藏的文本完全透明、动态，而且动态的范围很小
\beamertemplatetransparentcovereddynamic
% 使itemize环境中变成小球，这是一种视觉效果
\beamertemplateballitem
% 为所有已编号的部分设置一个章节目录，并且编号显示成小球
\beamertemplatenumberedballsectiontoc
% 将每一页的要素的要素名设成加粗字体
\beamertemplateboldpartpage

% item逐步显示时,使已经出现的item、正在显示的item、将要出现的item呈现不同颜色
\def\hilite<#1>{
 \temporal<#1>{\color{gray}}{\color{blue}}
    {\color{blue!25}}
}

\renewcommand{\today}{\number\year 年 \number\month 月 \number\day 日}

%五角星
\usepackage{MnSymbol}

%去除图表标题中的figure等
\usepackage{caption}
\captionsetup{labelformat=empty,labelsep=none}

\usepackage{tabu}
\usepackage{multirow}
%表格自动换行
\usepackage{tabularx} 

% 千分号
%\usepackage{textcomp}

%罗马数字
\makeatletter
\newcommand{\rmnum}[1]{\romannumeral #1}
\newcommand{\Rmnum}[1]{\expandafter\@slowromancap\romannumeral #1@}
\makeatother

%分栏
\usepackage{multicol}

%\usepackage{enumitem}
%\usepackage{enumerate}

%键盘
\usepackage{keystroke}

%心形
%\usepackage{fdsymbol}

%插入源代码
\usepackage{listings}
\lstset{
  language=perl,                  % 程序语言名称：TeX, Perl, R, sh, bash, Awk
  basicstyle=\normalsize\tt,      %\tt指monospace字体族，程序源代码使用此族字体表示更加美观
  numbers=left,                   % 行号位置(左侧)
  numberstyle=\small,             % 行号字体的字号
  stepnumber=1,                   % 行号的显示步长
  numbersep=5pt,                  % 行号与代码间距
  backgroundcolor=\color{white},  % 背景色;需要 \usepackage{color}
  showspaces=false,               % 不显示空格
  showstringspaces=false,         % 不显示代码字符串中的空格标记
  showtabs=false,                 % 不显示 TAB
  tabsize=4, 
  frame=shadowbox,                % 把代码用带有阴影的框圈起来
  captionpos=b,                   % 标题位置
  breaklines=true,                % 对过长的代码自动断行
  breakatwhitespace=false,        % 断行只在空格处
  extendedchars=false,            % 解决代码跨页时，章节标题，页眉等汉字不显示的问题
  %escapeinside={\%*}{*},         % 跳脱字符，添加注释，暂时离开 listings 
  %escapeinside=``,
  commentstyle=\color{red!50!green!50!blue!50}\tt,  %浅灰色的注释
  rulesepcolor=\color{red!20!green!20!blue!20},     %代码块边框为淡青色
  keywordstyle=\color{blue!70}\bfseries\tt,         %代码关键字的颜色为蓝色，粗体
  identifierstyle=\tt,
  stringstyle=\tt,                % 代码字符串的特殊格式
  keepspaces=true,
  breakindent=1em,
  %breakindent=22pt,
  %breakindent=4em,
  breakautoindent=true,
  flexiblecolumns=true,
  aboveskip=1em,                  %代码块边框
  xleftmargin=2em,
  xrightmargin=2em
}

%\setbeamercolor{alerted text}{fg=magenta}
\setbeamercolor{bgcolor}{fg=yellow,bg=cyan}
%\setbeamercolor{itemize/enumerate body}{fg=green}

\begin{document}

%\includeonlyframes{current}

\logo{\includegraphics[height=0.08\textwidth]{qr.png}}

% 在每个Section前都会加入的Frame
\AtBeginSection[]
{
  \begin{frame}<beamer>
    %\frametitle{Outline}
    \frametitle{教学提纲}
    \setcounter{tocdepth}{3}
    \begin{multicols}{2}
      \tableofcontents[currentsection,currentsubsection]
      %\tableofcontents[currentsection]
    \end{multicols}
  \end{frame}
}
% 在每个Subsection前都会加入的Frame
\AtBeginSubsection[]
{
  \begin{frame}<beamer>
%%\begin{frame}<handout:0>
%% handout:0 表示只在手稿中出现
    \frametitle{教学提纲}
    \setcounter{tocdepth}{3}
    \begin{multicols}{2}
    \tableofcontents[currentsection,currentsubsection]
    \end{multicols}
%% 显示在目录中加亮的当前章节
  \end{frame}
}

% 为当前幻灯片设置背景
%{
%\usebackgroundtemplate{
%\vbox to \paperheight{\vfil\hbox to
%\paperwidth{\hfil\includegraphics[width=2in]{tijmu_charcoal.png}\hfil}\vfil}
%}
\begin{frame}[plain]
  \begin{center}
    {\Huge 疯狂的实验\\}
    \vspace{1cm}
    {\LARGE 天津医科大学\\}
    %\vspace{0.2cm}
    {\LARGE 生物医学工程与技术学院\\}
    \vspace{1cm}
    {\large 2018-2019学年下学期（春）\\ 公共选修课}
  \end{center}
\end{frame}
%}



%\includeonlyframes{current}

\title[其他]{第五章\quad 疯狂的实验——其他}
\author[Yixf]{伊现富（Yi Xianfu）}
\institute[TIJMU]{天津医科大学（TIJMU）\\ 生物医学工程与技术学院}
\date{2017年4月}

\begin{frame}
  \titlepage
\end{frame}

\begin{frame}[plain,label=current]
  \frametitle{教学提纲}
  \setcounter{tocdepth}{3}
  \begin{multicols}{2}
    \tableofcontents
  \end{multicols}
\end{frame}


% 1-221,1970
\section{福克斯博士效应}
\begin{frame}
  \frametitle{福克斯博士效应 | 实验}
  \begin{block}{问题}
 是否可以用炫人耳目的演讲技巧欺骗一组专家，使他们注意不到内容的荒唐无稽？ 
  \end{block}
  \pause
  \begin{block}{实验}
    \begin{itemize}
      \item 福克斯：“将数学科学应用于人类行为的权威”（演员）
      \item 报告：《培训医生过程中所应用的数学博弈理论》
      \item 听众：经验丰富的教育学家
      \item 策划：报告充满了含混不清的言语、凭空捏造的词句和自相矛盾的结论，借助大量的幽默和参阅提示把这些胡言乱语的内容组织起来；在提问环节，使用高超巧妙的方式予以回答
    \end{itemize}
  \end{block}
\end{frame}

\begin{frame}
  \frametitle{福克斯博士效应 | 结果}
  \begin{block}{结果}
    \begin{itemize}
      \item 一小时后报告正文结束，大家积极提问。
      \item 全部10位听众都表示这次报告启迪了他们的思考。
      \item 9位认为福克斯材料整理清楚有趣、内容介绍妙趣横生并且穿插了丰富的阐释性事例。
      \item 播放报告录像，有一位居然认为他以前读过福克斯的论文。
      \item 得知福克斯的真实身份后，有几位听众想要询问更多的文献信息。
    \end{itemize}
  \end{block}
  \pause
  \begin{block}{福克斯博士效应}
    报告的风格完全可以演示贫乏的内容。
  \end{block}
\end{frame}

\begin{frame}
  \frametitle{福克斯博士效应 | 启示}
  \begin{block}{启示}
    \begin{itemize}
      \item 课程评估的说服力值得怀疑：学生们在问卷中写下的课程评价可能无非是他们自己的满意感和“学到了东西的错觉”。
      \item 尽管报告是一个一无所云的骗局，但它的风格唤起了听众对题目的兴趣。 $\Longrightarrow$ 建议人们创新方法提高学生的学习动力。
      \item 如果一个演员可以成为高明的老师，为什么不能成为出色的议员甚至优秀的总统呢？ $\Longrightarrow$ 1980年，里根成为美国总统！（他的演说风格高明而极具说服力，被媒体誉为“伟大的沟通者”，美国人心目中最伟大的总统之一。）
    \end{itemize}
  \end{block}
\end{frame}

\section{《圣经》}
% 2-22,1822
\begin{frame}
  \frametitle{圣经 | 鬣狗}
  \begin{block}{缘起}
    采石场的工人们发现了一个洞穴，里面堆满了老虎、牡鹿、熊、马、大象、犀牛、河马和鬣狗的骨头。
  \end{block}
  \vspace{-0.5em}
  \pause
  \begin{block}{解释与疑问}
    \begin{itemize}
      \item 解释：教士借助《圣经》故事——野兽们被“大洪水”冲进洞穴。
      \item 疑问：洞里没有沙石的踪迹。大型动物是怎样通过狭窄的洞口挤入洞中的呢？
    \end{itemize}
  \end{block}
  \vspace{-0.5em}
  \pause
  \begin{block}{猜测与实验}
    \begin{itemize}
      \item 观察：骨头上有被啃噬的痕迹，而且牙印与鬣狗牙齿完全吻合。
      \item 推断：鬣狗将觅得的腐尸拖进洞里慢慢享用。
      \item 实验：从非洲南部引进鬣狗，让它啃食骨头。
      \item 结果：不吃的部位都能在洞穴里找到，吃掉的部位证实洞穴里没有的。实验用的骨头和洞穴里的骨头的断裂状况完全一样。
    \end{itemize}
  \end{block}
\end{frame}

% 1-292,1998
\begin{frame}
  \frametitle{圣经 | 耶利哥的扬声器}
  \begin{block}{耶利哥的扬声器}
    \begin{itemize}
      \item 7个祭祀在约柜前吹响他们的羊角，使得耶利哥的城墙倾陷——《圣经\textbullet 约书亚记》
      \item 美国教育节目《学习频道》，加利福尼亚的怀勒实验室，怀勒WAS 3000扬声器（相当于10000个扬声喇叭）
      \item 结果：经过6分钟的噪声作用，砂浆真的开始破碎，小墙分崩离析
      \item 结论：（众所周知的事实）声音真的可以产生破坏
      \item 事实：耶利哥所属的迦南城根本没有设防……
    \end{itemize}
  \end{block}
\end{frame}

% 
\section{}

% 
\section{}

% 
\section{}


\section{史海撷华}
% 2-23,1874
\begin{frame}
  \frametitle{史海撷华 | 创伤弹道学}
  \begin{block}{疑问}
    枪击究竟为何能够致死？
  \end{block}
  \pause
  \begin{block}{猜测}
    \begin{enumerate}
      \item “碎片”说：极高的温度使子弹熔化，导致许多碎片从子弹上剥落下来。
      \item “离心力”说：旋转的子弹带有离心力，折断了皮肉。
      \item “压力”说：子弹侵入肌肉以及柔软部位后产生了一种压力。
    \end{enumerate}
  \end{block}
  \pause
  \begin{block}{观察}
    \begin{itemize}
      \item 子弹飞出体外的地方、破掉的皮肉并没有扭转出涡旋式的形状——“离心力”的说法并不准确。
      \item 如果子弹没有击中骨骼就不会留下什么碎片——“碎片”的说法也不可靠。
    \end{itemize}
  \end{block}
\end{frame}

\begin{frame}
  \frametitle{史海撷华 | 创伤弹道学}
  \begin{block}{实验}
    \begin{itemize}
      \item 主持：科赫尔（首位获得诺贝尔医学奖的外科医生，1909年）
      \item 武器：口径为10.4、11毫米的“维特尔利”、“夏斯波”
      \item 靶子：杉木板、书、装满沙子的猪膀胱、塞满土豆泥的人类颅骨、用布裹住的尸体、石板、铁皮罐、玻璃板、肝脏……
    \end{itemize}
  \end{block}
  \vspace{-0.3em}
  \pause
  \begin{block}{结果}
     子弹打在空的颅骨上——两个洞；打在用土豆泥填充的颅骨上——开花。
  \end{block}
  \vspace{-0.3em}
  \pause
  \begin{block}{启示}
    \begin{itemize}
      \item 枪击造成的流体静力冲击才是摧毁人体组织的元凶。
      \item 为“鲁宾弹”的设计提供理论基础；开创“创伤弹道学”研究先河。
      \item 战争的目的不是尽量消灭更多人的生命，而是“把一个拥有较强战斗能力的敌对者变成一个需要照料的病人”。
    \end{itemize}
  \end{block}
\end{frame}

% 1-193,1966
\begin{frame}
  \frametitle{史海撷华 | 搭车技巧 | 弱不禁风}
  \begin{block}{初衷}
    研究汽车司机对待拦车搭乘者的态度。
  \end{block}
  \pause
  \begin{block}{实验}
    男生，用4天时间站在一条4车道的马路边尝试搭车，有时膝盖裹着绷带并拄着拐杖，有时则没有。
  \end{block}
  \pause
  \begin{block}{结果}
    绑绷带拄拐杖时获得了至少双重的搭车机会 $\Longrightarrow$ 搭车技巧之一：要弱不禁风！
  \end{block}
\end{frame}

% 1-231,1971
\begin{frame}
  \frametitle{史海撷华 | 搭车技巧 | 是个女的}
  \begin{block}{《搭乘的成功率》}
    \begin{itemize}
      \item 研究男女在穿着不同服装时，站在路边等待搭乘的成功率。
      \item 研究不同人数和搭配对司机态度的影响。
    \end{itemize}
  \end{block}
  \pause
  \begin{block}{结论}
    \begin{itemize}
      \item 衣着邋遢的比衣着整洁干净的难以搭车。
      \item 男人比女人更难搭车。
      \item 2个男人是最难搭车的。
      \item 单独一个男人搭车的成功率根一对情侣相差无几。
      \item 成功率最高的是2个女人在一起时。
    \end{itemize}
  \end{block}
\end{frame}

% 1-231,1971
\begin{frame}
  \frametitle{史海撷华 | 搭车技巧 | 更多技巧}
  \begin{block}{目光交流}
    \begin{itemize}
      \item 当被注视时，加州帕罗奥多的600辆车中，有40辆车停了下来。
      \item 如果司机和搭乘者之间没有目光交流，仅有18辆车会停下来。
    \end{itemize}
  \end{block}
  \pause
  \begin{block}{丰胸}
在西雅图进行的一次实验中，带着垫厚的胸罩（加上5厘米的外延）的女士获得了与不这样做时相比双倍的搭乘机会。 
  \end{block}
\end{frame}

% 1-254,1976 
\begin{frame}
  \frametitle{史海撷华 | 剃须刀教学法}
  \begin{block}{实验室结论}
    \begin{itemize}
      \item “对作为客体的人的稳定的刺激特征”（胡子）会在对人性格的判断上产生深入的影响。
      \item 脸上的一个细小变化，就会让观察者获得不同的印象。
    \end{itemize}
  \end{block}
  \vspace{-0.5em}
  \pause
  \begin{block}{疑问}
      实验室研究中得出的上述结论在现实中是否成立呢？
      %《胡子的魔术——关于长胡子的高校教员对学生所产生的影响的研究》
  \end{block}
  \vspace{-0.5em}
  \pause
  \begin{block}{实验}
    \begin{itemize}
      \item 不留胡子（2学期）——留胡子（1学期）——不留胡子（1学期）
      \item 学期开始见面10分钟后，填写评价问卷
    \end{itemize}
  \end{block}
  \vspace{-0.5em}
  \pause
  \begin{block}{结果}
留胡子：不够有目标、不够专心、不够友善、不够坚强、不够理智、不够聪明……；看起来不拘束、比较随意……
  \end{block}
\end{frame}

% 
\begin{frame}
  \frametitle{史海撷华 | }
\end{frame}

% 
\begin{frame}
  \frametitle{史海撷华 | }
\end{frame}

% 
\begin{frame}
  \frametitle{史海撷华 | }
\end{frame}

% 
\begin{frame}
  \frametitle{史海撷华 | }
\end{frame}

% 
\begin{frame}
  \frametitle{史海撷华 | }
\end{frame}

% 
\begin{frame}
  \frametitle{史海撷华 | }
\end{frame}

% 
\begin{frame}
  \frametitle{史海撷华 | }
\end{frame}

% 
\begin{frame}
  \frametitle{史海撷华 | }
\end{frame}



\section*{Acknowledgements}
\begin{frame}
  \frametitle{Powered by}
  \begin{center}
    \includegraphics[width=9cm]{power.png}
  \end{center}
\end{frame}

\end{document}

